% Topic 5: Inverses and Determinants for nxn Matrices
% Part of the Matrix Fundamentals section

\subsection{Inverses and Determinants for $n \times n$ Matrices}

While the inverse of a $2 \times 2$ matrix has a simple formula, extending this to $3 \times 3$ matrices (and larger) requires a more systematic process using \textbf{minors} and \textbf{cofactors}.

\subsubsection*{Minors and Cofactors}

For any element $a_{ij}$ in a square matrix $A$:
\begin{itemize}[leftmargin=*]
  \item The \textbf{minor}, $M_{ij}$, is the determinant of the smaller matrix that remains after deleting row $i$ and column $j$.
  \item The \textbf{cofactor}, $C_{ij}$, is the minor multiplied by a positive or negative sign, based on its position:
  \[
  C_{ij} = (-1)^{i+j} M_{ij}
  \]
  This creates a ``checkerboard'' pattern of signs. For a $3 \times 3$ matrix:
  \[
  \begin{bmatrix} + & - & + \\ - & + & - \\ + & - & + \end{bmatrix}
  \]
\end{itemize}

\subsubsection*{Evaluating a $3 \times 3$ Determinant}

The determinant of a $3 \times 3$ matrix can be evaluated by \textbf{cofactor expansion} along any row or column. You multiply each element in that row (or column) by its corresponding cofactor and add them up.

Expanding along the first row:
\[
\det(A) = a_{11}C_{11} + a_{12}C_{12} + a_{13}C_{13}
\]
Using the $2 \times 2$ determinant formula for the minors, this expands to:
\[
\det(A) = a_{11}(a_{22}a_{33} - a_{23}a_{32}) - a_{12}(a_{21}a_{33} - a_{23}a_{31}) + a_{13}(a_{21}a_{32} - a_{22}a_{31})
\]

\subsubsection*{The Inverse of an $n \times n$ Matrix}

The formula for the inverse of any non-singular $n \times n$ matrix is:
\[
A^{-1} = \frac{1}{\det(A)} \text{adj}(A)
\]
Where $\text{adj}(A)$, called the \textbf{adjugate matrix}, is the \emph{transpose} of the matrix of cofactors.

\subsubsection*{Geometric Meaning: Determinants as Area Engines}

Determinants are not merely computational recipes for finding inverses; they possess a direct geometric interpretation. In two dimensions, the determinant of a $2 \times 2$ matrix gives the signed area of the parallelogram formed by its column vectors.

For any triangle with three vertices $P_1(x_1, y_1)$, $P_2(x_2, y_2)$, and $P_3(x_3, y_3)$ in the Cartesian plane, its area is given by the $3 \times 3$ determinant:
\[
\text{Area}(\triangle P_1 P_2 P_3) = \frac{1}{2} \left| \det \begin{bmatrix} x_1 & y_1 & 1 \\ x_2 & y_2 & 1 \\ x_3 & y_3 & 1 \end{bmatrix} \right|.
\]
We can understand why this formula works using elementary row operations. Subtracting the first row from the second and third rows (which geometrically translates the triangle so that vertex $P_1$ moves to the origin $(0,0)$):
\[
\det \begin{bmatrix} x_1 & y_1 & 1 \\ x_2 & y_2 & 1 \\ x_3 & y_3 & 1 \end{bmatrix}
= \det \begin{bmatrix} x_1 & y_1 & 1 \\ x_2 - x_1 & y_2 - y_1 & 0 \\ x_3 - x_1 & y_3 - y_1 & 0 \end{bmatrix}.
\]
Expanding along the third column, the $3 \times 3$ determinant collapses immediately into a $2 \times 2$ determinant:
\[
\text{Area} = \frac{1}{2} \left| (x_2 - x_1)(y_3 - y_1) - (x_3 - x_1)(y_2 - y_1) \right|.
\]
If one vertex is already at the origin $O(0,0)$, the formula simplifies directly to:
\[
\text{Area}(\triangle OP_2 P_3) = \frac{1}{2} |x_2 y_3 - x_3 y_2| = \frac{1}{2} \left| \det \begin{bmatrix} x_2 & y_2 \\ x_3 & y_3 \end{bmatrix} \right|.
\]

\begin{remark}[Collinearity Criterion]
Three points $P_1, P_2, P_3$ are \textbf{collinear} (lie on the same straight line) if and only if the area of the triangle they form is zero:
\[
\det \begin{bmatrix} x_1 & y_1 & 1 \\ x_2 & y_2 & 1 \\ x_3 & y_3 & 1 \end{bmatrix} = 0.
\]
\end{remark}

\begin{remark}[Concurrence Criterion for Three Lines]
Historically, the 1965 NSW First Level syllabus Section~4(a) explicitly paired the triangle area formula with the determinant condition for the concurrence of three straight lines.

Three distinct lines in the Cartesian plane:
\[
L_1: a_1 x + b_1 y + c_1 = 0, \qquad
L_2: a_2 x + b_2 y + c_2 = 0, \qquad
L_3: a_3 x + b_3 y + c_3 = 0
\]
are \textbf{concurrent} (pass through a common single point) if and only if the determinant of their coefficients vanishes:
\[
\det \begin{bmatrix} a_1 & b_1 & c_1 \\ a_2 & b_2 & c_2 \\ a_3 & b_3 & c_3 \end{bmatrix} = 0.
\]
This highlights a remarkable projective duality in coordinate geometry:
\begin{itemize}[leftmargin=*]
  \item \textbf{Three points are collinear} $\iff$ the coordinate matrix has determinant zero.
  \item \textbf{Three lines are concurrent} $\iff$ the line coefficient matrix has determinant zero.
\end{itemize}
\end{remark}

\bigskip

% ── Worked Example: 3x3 Inverse ──
\begin{problem}[Inverting a $3 \times 3$ Matrix]
Let $A = \begin{bmatrix} 1 & 2 & -1 \\ 0 & 3 & 1 \\ 2 & 0 & -2 \end{bmatrix}$.

\begin{enumerate}[label=\textbf{(\alph*)}]
  \item Evaluate $\det(A)$ by expanding along the first row.
  \item Find the matrix of cofactors, $C$.
  \item Find the adjugate matrix, $\text{adj}(A)$.
  \item State the inverse matrix, $A^{-1}$.
\end{enumerate}
\end{problem}

\begin{solution}
\textbf{(a) Determinant} \\
Expanding along the first row:
\begin{align*}
\det(A) &= 1 \begin{vmatrix} 3 & 1 \\ 0 & -2 \end{vmatrix} - 2 \begin{vmatrix} 0 & 1 \\ 2 & -2 \end{vmatrix} + (-1) \begin{vmatrix} 0 & 3 \\ 2 & 0 \end{vmatrix} \\[6pt]
&= 1( (3)(-2) - (1)(0) ) - 2( (0)(-2) - (1)(2) ) - 1( (0)(0) - (3)(2) ) \\
&= 1( -6 ) - 2( -2 ) - 1( -6 ) \\
&= -6 + 4 + 6 = 4
\end{align*}
Since $\det(A) = 4 \neq 0$, the inverse exists.

\medskip
\textbf{(b) Matrix of Cofactors} \\
We calculate the cofactor for all 9 elements. Remember to apply the checkerboard signs!
\begin{align*}
C_{11} &= +( (3)(-2) - (1)(0) ) = -6 \\
C_{12} &= -( (0)(-2) - (1)(2) ) = -(-2) = 2 \\
C_{13} &= +( (0)(0) - (3)(2) ) = -6 \\[6pt]
C_{21} &= -( (2)(-2) - (-1)(0) ) = -(-4) = 4 \\
C_{22} &= +( (1)(-2) - (-1)(2) ) = -2 + 2 = 0 \\
C_{23} &= -( (1)(0) - (2)(2) ) = -(-4) = 4 \\[6pt]
C_{31} &= +( (2)(1) - (-1)(3) ) = 2 + 3 = 5 \\
C_{32} &= -( (1)(1) - (-1)(0) ) = -(1) = -1 \\
C_{33} &= +( (1)(3) - (2)(0) ) = 3
\end{align*}
The cofactor matrix is:
\[
C = \begin{bmatrix} -6 & 2 & -6 \\ 4 & 0 & 4 \\ 5 & -1 & 3 \end{bmatrix}
\]

\medskip
\textbf{(c) Adjugate Matrix} \\
The adjugate is the transpose of the cofactor matrix ($C^T$):
\[
\text{adj}(A) = \begin{bmatrix} -6 & 4 & 5 \\ 2 & 0 & -1 \\ -6 & 4 & 3 \end{bmatrix}
\]

\medskip
\textbf{(d) Inverse Matrix} \\
Using the formula $A^{-1} = \frac{1}{\det(A)}\text{adj}(A)$:
\[
A^{-1} = \frac{1}{4} \begin{bmatrix} -6 & 4 & 5 \\ 2 & 0 & -1 \\ -6 & 4 & 3 \end{bmatrix}
\]
\end{solution}

\begin{takeaways}
\item \textbf{Row Choice for Determinants:} You can expand along \emph{any} row or column to find the determinant. For speed, always choose the row or column containing the most zeros.
\item \textbf{Don't Forget the Transpose:} A common error is writing the cofactor matrix directly into the inverse formula. You must take the transpose of the cofactor matrix to get the adjugate.
\item \textbf{Sign Errors:} Calculating a $3 \times 3$ inverse manually involves many small multiplications and subtractions. Write out your working explicitly to avoid dropping minus signs.
\item \textbf{Geometric Area Engine:} The determinant of a $3 \times 3$ matrix with a column of 1s computes the area of any triangle in the plane, bypassing the need for distance and perpendicular height formulas.
\end{takeaways}
